One of the most fundamental questions in statistics is deceptively simple: given a dataset,
does it come from a distribution we have in mind, or not? This is the goodness-of-fit
testing problem, and it has been studied for over a century. Classical tests such as
Kolmogorov--Smirnov or Pearson's $\chi^2$ were designed for simple, low-dimensional
settings. But as modern scientific applications generate data of increasing complexity and
dimension, a natural question arises: how many observations do we actually need to answer
this question reliably, and can we do so efficiently?

This project addresses these questions in a nonparametric minimax framework, where one
seeks to characterize the fundamental statistical limits of the problem, that is, the
minimum number of samples required by \emph{any} test, regardless of its construction.
The setting we focus on is the \textbf{Gaussian sequence model} with mean vector constrained
to a Sobolev ellipsoid, a canonical model in nonparametric statistics that captures the
essential difficulty of signal detection in function spaces. In this model, the unknown
signal is represented as an infinite sequence of Gaussian observations, and the smoothness
parameter $s$ of the Sobolev ellipsoid controls how fast the signal's frequency components
decay. This makes it possible to study, in a precise and clean way, how the regularity of
the signal affects the difficulty of testing.

Two approaches to this problem are developed in this project. The first, covered in
Chapter~3, is the classical approach based on $\chi^2$-type tests, pioneered by \cite{Ing2012}. These tests aggregate a weighted sum of squared
observations and reject the null hypothesis when this sum is large. The minimax theory
for such tests is well understood, and they are known to be optimal in the Gaussian
sequence model. The second approach, covered in Chapter~4, is more recent and comes from
a surprising direction: machine learning. The \textbf{Classifier-Accuracy Test} (CAT),
introduced by \cite{gerberclas}, replaces the classical test statistic by
a binary classifier trained to distinguish samples from the null and alternative
distributions. Rather than computing a closed-form statistic, the classifier implicitly
learns a \emph{separating set} (i.e. a region of the sample space where the two
distributions differ most) and the test is based on how many observations fall in
this region.

At first glance, restricting the classifier to a binary output might seem wasteful in
view of the Neyman--Pearson lemma. Yet \cite{gerberclas} shows that a well-chosen binary
separating set already achieves minimax-optimal sample complexity, making classifier-based
testing a theoretically grounded procedure rather than a mere heuristic.

The project is organized as follows. \Cref{chap2} introduces the general minimax framework
for hypothesis testing, including the Bayesian and minimax approaches and the connection
between them. \Cref{chap3} specializes to the Gaussian sequence model and develops the
classical theory of minimax distinguishability and $\chi^2$-type tests. \Cref{chap4}
introduces the CAT framework, establishes the key lemmas needed to analyze the performance
of the separating set, and proves the minimax lower bound on the sample complexity.

\section*{Acknowledgement}
I am deeply grateful to Prof.~Sven Wang for offering me the opportunity to undertake
this project under his supervision. I would also like to extend my heartfelt thanks to
Duc Hoang for the considerable time he devoted to guiding me, for his thoughtful and
constructive feedback, and for the freedom he gave me in shaping the content of this
work.